\documentclass[11pt]{article}
\usepackage[letterpaper,margin=1in]{geometry}
\usepackage{amsmath,amssymb,amsthm}
\usepackage{booktabs,tabularx,array}
\usepackage[colorlinks=true,linkcolor=black,citecolor=black,urlcolor=blue]{hyperref}
\hypersetup{
  pdftitle={Counterexamples to a circular coloring conjecture for the Cartesian sum},
  pdfauthor={Alex Chan}}

\newtheorem{theorem}{Theorem}
\newtheorem{lemma}[theorem]{Lemma}
\newtheorem{corollary}[theorem]{Corollary}
\newcommand{\out}{N^{+}}
\newcommand{\Z}{\mathbb Z}
\newcommand{\cc}{\chi_c}

\title{Counterexamples to a circular coloring conjecture for the Cartesian sum}
\author{Alex Chan\\[2pt]
  \small Harvard University and National Bureau of Economic Research\\[-1pt]
  \small \texttt{achan@hbs.edu}}
\date{\today}

\begin{document}
\maketitle

\begin{abstract}
Liu and Zhu conjectured that the circular chromatic number of the Cartesian
sum, also called the disjunctive or OR product, is at most
$\max\{\cc(G)\chi(H),\chi(G)\cc(H)\}$.  We give an infinite family of
counterexamples with both factors circular complete graphs.  The smallest
instance here is $K_{5/2}\oplus K_{11/4}$, a $34$-regular graph on $55$
vertices with $33/4<\chi_c\leq17/2$.
The proof reduces a hypothetical coloring to two rigid arrangements of
intervals on a circle and excludes both arrangements.
\end{abstract}

\section{Statement}
For integers $p\geq 2q>0$, let $K_{p/q}$ have vertex set $\Z_p$, with
distinct $x,y$ adjacent when their cyclic distance is at least $q$.  The
Cartesian sum $G\oplus H$ has vertex set $V(G)\times V(H)$; two pairs are
adjacent when at least one pair of corresponding coordinates is adjacent.
A $(P,Q)$-coloring is a homomorphism to $K_{P/Q}$.  Liu and Zhu
\cite{LiuZhu2008} conjectured
\begin{equation}\label{eq:conj}
 \cc(G\oplus H)\leq
 \max\{\cc(G)\chi(H),\chi(G)\cc(H)\}.
\end{equation}
Their theorem concerning the ordinary chromatic number is a
distinct statement.\footnote{A Lean~4 certificate for the $55$-vertex
instance is available at
\href{https://github.com/alexgocardinal/liu-zhu-circular-coloring}{the accompanying GitHub repository}
(commit \texttt{a4b5b22}). It verifies the nonexistence of a
$(33,4)$-coloring and the explicit $(17,2)$-coloring below.
The certificate uses Lean and Mathlib version~4.35.0; the obstruction
depends only on the standard logical axioms
\texttt{propext}, \texttt{Classical.choice}, and \texttt{Quot.sound}.
The infinite-family theorem and its circular-chromatic interpretation
are proved in the text and are outside the formalization's scope.}

\begin{theorem}\label{thm:main}
Let $q\geq4$, let $p\in\{3q-2,3q-1\}$, and assume
$\gcd(p,q)=\gcd(q,3)=1$.  Then
\[
 \cc\bigl(K_{5/2}\oplus K_{p/q}\bigr)>
 \frac{3p}{q}
 =\max\{\cc(K_{5/2})\chi(K_{p/q}),
          \chi(K_{5/2})\cc(K_{p/q})\}.
\]
These graphs have $5p$ vertices and degree $5p-6q+3$.
\end{theorem}

In particular, $p=3q-1$ is admissible for every $q\geq4$ not divisible
by $3$.  This gives infinitely many orders.  The $p=3q-2$ choice is
admissible whenever $q$ is also odd.

\begin{corollary}\label{cor:55}
The graph $X=K_{5/2}\oplus K_{11/4}$ has $55$ vertices, $935$ edges, and
degree $34$, and satisfies $33/4<\cc(X)\leq17/2$ and $\chi(X)=9$.
\end{corollary}
\begin{proof}
Apply Theorem~\ref{thm:main} with $(p,q)=(11,4)$.  The coloring
$(a,b)\mapsto3\lfloor a/2\rfloor+\lfloor b/4\rfloor$ takes values in
$\{0,\ldots,8\}$ and is proper: each color class is contained in a
product of two consecutive blocks, independent in the respective
factors.  The explicit $(17,2)$-coloring in Table~\ref{tab:witness}
gives the other circular bound.  Since $\chi(X)=\lceil\cc(X)\rceil$, the strict lower bound gives
$\chi(X)\geq9$.  The degree follows by counting nonneighbors: a vertex
has $(2\cdot2-1)(2\cdot4-1)=21$ possible coordinatewise nonneighbors,
including itself, out of $55$ vertices.  Hence its degree is $55-21=34$.
\end{proof}

\begin{table}[ht]
\centering\small
\begin{tabular}{c*{11}{r}}
\toprule
 $a\backslash b$ & 0&1&2&3&4&5&6&7&8&9&10\\\midrule
 0&0&0&11&11&11&11&6&6&6&0&0\\
 1&16&15&15&10&10&10&5&5&5&5&16\\
 2&15&15&14&14&9&9&9&4&4&4&3\\
 3&2&13&13&13&13&8&8&8&3&3&3\\
 4&1&1&12&12&12&7&7&7&7&1&1
\\
\bottomrule
\end{tabular}
\caption{A $(17,2)$-coloring of $X$: the entry in row $a$, column $b$
is the color of $(a,b)\in\Z_5\times\Z_{11}$.  Every edge has target
cyclic color distance at least $2$.}\label{tab:witness}
\end{table}

\begin{corollary}[identical factors]\label{cor:diagonal}
The graph $K_{11/4}\oplus K_{11/4}$ violates \eqref{eq:conj}.
It has $121$ vertices and $4356$ edges.
\end{corollary}
\begin{proof}
The map $a\mapsto\lfloor11a/5\rfloor$ sends $K_{5/2}$
homomorphically into $K_{11/4}$.  Taking its product with the identity
on $K_{11/4}$ sends $X$ into the stated graph, whose circular chromatic
number is therefore greater than $33/4$.  The proposed right side for
these two identical factors is $33/4$.  The edge count follows from
degree $121-(2\cdot4-1)^2=72$.
\end{proof}

\section{A saturated interval system}
We prove the theorem.  Put $G=K_{5/2}$, $H=K_{p/q}$, and $P=3p$.
The independent sets of $G$ are singletons and consecutive pairs in a
geometric $5$-cycle (whose edges are the nonedges of $G$).

Every independent set of $H$ is contained in a cyclic arc of $q$
consecutive residues.  Indeed, translate one of its elements to $0$.
All other elements lie in $[0,q-1]\cup[p-q+1,p-1]$.  If $u$ and $-v$
are the extreme elements on the two sides and $u+v\geq q$, then their
compatibility forces $p-(u+v)\leq q-1$; since $u+v\leq2q-2$, this gives
$p\leq3q-3$, contrary to the hypothesis.  Thus $u+v\leq q-1$.
Consequently $\alpha(H)=q$, and its maximum independent sets are
precisely its cyclic $q$-arcs.  The identity $(p,q)$-coloring and the
fractional lower bound $|V(H)|/\alpha(H)$ show $\cc(H)=p/q$.
Three consecutive $q$-arcs give a $3$-coloring, while $p/q>2$ shows
$\chi(H)=3$.  Similarly $\cc(G)=5/2$ and $\chi(G)=3$.  As $p/q\geq5/2$,
the right side of \eqref{eq:conj} is $3p/q$.  Our coprimality assumptions
give $\gcd(P,q)=1$.  It suffices to show that there is no $(P,q)$-coloring
of $G\oplus H$: for a finite graph, $\cc(F)\leq P/q$ if and only if
$F\to K_{P/q}$ \cite{LiuZhu2008}.

Suppose such a coloring $c$ exists.  Replace each color $c(a,b)$ by its
interval $I(a,b)$ of $q$ consecutive slots in $\Z_P$.  Intervals
corresponding to adjacent vertices are disjoint.  Write
$U_b=\bigcup_{a\in\Z_5}I(a,b)$.  We claim $|U_b|\geq3q$.
If $U_b$ is not the whole circle, cut it at an unoccupied slot, and let
$J$ be the intersection graph of its five line intervals.  Since $G$
is a subgraph of the complement of $J$, the interval-graph coloring
property gives
\[
 3=\chi(G)\leq\chi(\overline J)=\alpha(J)
 \leq\left\lfloor\frac{|U_b|}{q}\right\rfloor.
\]
The claim is immediate if $U_b=\Z_P$.  At a target slot, the active
second coordinates form an independent set of $H$, so at most $q$ fibers
are active.  Therefore
\begin{equation}\label{eq:occupancy}
 3pq\leq\sum_{b\in\Z_p}|U_b|\leq Pq=3pq.
\end{equation}
Equality gives $|U_b|=3q<P$ for all $b$ and exactly $q$ active fibers
at each slot.  They form a cyclic $q$-arc in $H$.

Each occupied run of $U_b$ has length divisible by $q$.  For if its run
lengths are $\ell_1,\ldots,\ell_k$, disjoint length-$q$ intervals fit
at most $\sum_i\lfloor\ell_i/q\rfloor$ times, whereas the preceding
interval argument requires at least three.  Hence
\[
 3\leq\sum_i\left\lfloor\ell_i/q\right\rfloor
 \leq\sum_i\ell_i/q=3.
\]
Partition the runs into $q$-blocks; there are three for each fiber.
Let $x_t$ be the number of blocks beginning at slot $t\in\Z_P$.
Every slot is covered by exactly $q$ blocks, giving
$\sum_{j=0}^{q-1}x_{t-j}=q$.  Subtracting consecutive equations gives
$x_t=x_{t-q}$.  Since $\gcd(P,q)=1$, we have $x_t=1$ for every $t$.
Let $b_t$ denote the second coordinate of the unique block beginning
at $t$.  This is a cyclic word of length $P$; each label occurs three
times, and the labels of every $q$ consecutive blocks are distinct
and form a cyclic $q$-arc in $H$.

\section{The two possible block words}
\begin{lemma}\label{lem:phase}
Up to reversing the cyclic order of $H$, exactly one of the following
forms can arise from the preceding block word:
\begin{enumerate}
\item Each $U_b$ is the union of three isolated $q$-runs.
\item Each $U_b$ is one $3q$-run, and the $p$ run starts occur every
three target slots.
\end{enumerate}
\end{lemma}
\begin{proof}
At time $t$ consider the ordered queue $(b_t,\ldots,b_{t+q-1})$.
Its underlying set is a cyclic $q$-arc.  Two successive queues have
$q-1$ common labels.  Since $p>2q$, their arcs either agree or differ
by shifting one endpoint one place.  Let $D_t$ be the net forward
displacement of the arc up to time $t$, lifting its starting label to
the integers.  Give each queue entry its start-time residue modulo $q$.
The departing and arriving entries have the same residue.  If $\tau_k$
is the residue initially at spatial position $k$ in the arc, then the
residue at position $k$ at time $t$ is $\tau_{k+D_t}$.

After $L=P=3p$ steps, the ordered queue returns to itself, while all
start-time residues have increased by $L$ modulo $q$.  The arc has
completed an integer number $W$ of turns, so $D_L=pW$ and
\begin{equation}\label{eq:closure}
 \tau_{k+D_L}=\tau_k+L\pmod q.
\end{equation}
The $\tau_k$ are a permutation of $\Z_q$.  If $g=\gcd(D_L,q)$,
iterating \eqref{eq:closure} $q/g$ times gives
$(q/g)L\equiv0\pmod q$; since $\gcd(L,q)=1$, we must have $g=1$.
Reverse the labeling of
$H$ if necessary so that $D_L>0$, hence $W\geq1$.
Equation \eqref{eq:closure} also forces
$\tau_k=\tau_0+ak\pmod q$ for some $1\leq a<q$.

Set $z_t=t-\tau_0-aD_t\pmod q$.  A hold advances $z$ by one.
A forward move is possible only at $z=0$ and sends $z$ to $1-a$;
a backward move is possible only at $z=-a$ and sends $z$ to $1$.
This directed phase graph consists of the ordinary $q$-cycle of holds
and two chords.  Its simple directed cycles are: $q$ holds; one forward
move and $a-1$ holds; or one backward move and $q-a-1$ holds.
In particular, if their multiplicities in a closed-walk decomposition
are $N_0,N_+,N_-$, then
\[
 L=qN_0+aN_++(q-a)N_-,\qquad D_L=N_+-N_-.
\]
Substituting $D_L=pW$ yields
\begin{equation}\label{eq:phase-arithmetic}
 3p=apW+q(N_0+N_-).
\end{equation}
Thus $1\leq aW\leq3$ and, by $\gcd(p,q)=1$,
$aW\equiv3\pmod q$.  As $q\geq4$, necessarily $aW=3$ and
$N_0=N_-=0$.  Either $(a,W)=(1,3)$, a forward step every slot, or
$(a,W)=(3,1)$, a forward step every three slots.  In the first case
each label traverses the arc three times, occupying three isolated
$q$-runs.  In the second it traverses the arc once, occupying a
single $3q$-run, with new runs starting every three slots.
\end{proof}

\section{Excluding both forms}
\begin{lemma}[cyclic interval expansion]\label{lem:expansion}
Let $A$ be a nonempty subset of $\Z_L$, and let $1\leq r\leq L$.
The union of the length-$r$ cyclic intervals starting at the elements
of $A$ has size at least
\[
  \min\{L,|A|+r-1\}.
\]
\end{lemma}
\begin{proof}
Set $B_0=A$ and $B_{j+1}=B_j\cup(B_j+1)$. A nonempty proper
subset of $\Z_L$ cannot be invariant under translation by one.
Consequently each step increases the cardinality by at least one
until the whole circle is occupied. The required union is $B_{r-1}$.
\end{proof}

Consider the first form of Lemma~\ref{lem:phase}. Each vertex
interval must coincide with one of its fiber's three isolated
$q$-blocks. Fix $a\in\Z_5$. The $p$ vertices $(a,b)$, with
$b\in\Z_p$, have distinct interval starts: a slot is the start of
exactly one block, and that block belongs to exactly one fiber.
Thus, writing
\[
  V_a=\bigcup_{b\in\Z_p} I(a,b),
\]
Lemma~\ref{lem:expansion} gives $|V_a|\geq p+q-1$, since
$p+q-1<P=3p$.

At any target slot, the active first coordinates form an independent
set of $K_{5/2}$, so there are at most two. Double counting gives
\[
  5(p+q-1)\leq\sum_{a\in\Z_5}|V_a|\leq2P=6p.
\]
It follows that $5(q-1)\leq p$. This contradicts
$p\leq3q-1<5(q-1)$ for $q\geq4$, excluding the first form.

\begin{lemma}[three shifted segments]\label{lem:segments}
Suppose three proper assignments of length-$q$ intervals to the
vertices of $K_{5/2}$ are contained in respective segments of length
$3q$, starting at $0,d,2d$, where $0<d<q$.  If intervals associated
with adjacent vertices in different assignments are also disjoint,
then $2d=q$.
\end{lemma}
\begin{proof}
For any one assignment, orient each edge of $K_{5/2}$ toward the
larger interval start.  An odd cycle has a directed two-edge path.
Its two consecutive start differences are at least $q$, while all
starts lie in an interval of length $2q$.  The three starts must
therefore be $0,q,2q$.  Label those vertices $0,2,4$ in the
complementary $5$-cycle.  The other starts have the form $x,q+y$,
where $0\leq x\leq y\leq q$.

Normalize the first assignment to have anchors $0,2,4$ at
$0,q,2q$.  The first anchor of the second assignment, at $d$,
overlaps the first anchors $0$ and $2$, so its label must be $1$,
the unique label nonadjacent to both.  Its middle anchor at $q+d$
overlaps anchors $2$ and $4$ of the first assignment and must
have label $3$.  Repeating this reasoning, the third assignment's
first anchor at $2d$ has label $2$.  It must be disjoint from the
first assignment's anchors $0$ at $0$ and $4$ at $2q$.
These two requirements give $2d\geq q$ and $2d\leq q$.
\end{proof}

In the second form of Lemma~\ref{lem:phase}, three successive fiber
runs start at $0,3,6$ after translation.  Their $3q$-segments fit
in a line because $3q+6<3p$.  Intervals for adjacent first
coordinates in different fibers are disjoint by the definition of
the Cartesian sum.  Lemma~\ref{lem:segments}, with $d=3$, forces
$q=6$, contradicting $\gcd(q,3)=1$.  Finally, each vertex has
$3(2q-1)$ coordinatewise nonneighbors, including itself, among the
$5p$ vertices; hence its degree is $5p-6q+3$.  This finishes the proof
of Theorem~\ref{thm:main}. \qed

\begin{thebibliography}{1}
\bibitem{LiuZhu2008}
D. D.-F. Liu and X. Zhu, Coloring the Cartesian sum of graphs,
\emph{Discrete Mathematics} \textbf{308} (2008), 5928--5936.
\href{https://doi.org/10.1016/j.disc.2007.11.005}{doi:10.1016/j.disc.2007.11.005}.
\end{thebibliography}
\end{document}
